\documentclass[11pt,a4paper]{article}

\usepackage[utf8]{inputenc}
\usepackage[T1]{fontenc}
\usepackage[margin=2.3cm]{geometry}
\usepackage{amsmath,amssymb}
\usepackage{booktabs}
\usepackage{longtable}
\usepackage{array}
\usepackage{xcolor}
\usepackage{graphicx}
\usepackage{caption}
\usepackage{subcaption}
\usepackage{tikz}
\usepackage{enumitem}
\usepackage{titlesec}
\usepackage{hyperref}
\usepackage[expansion=false]{microtype}
\usepackage{listings}

\hypersetup{
    colorlinks=true,
    linkcolor=blue!60!black,
    citecolor=blue!60!black,
    urlcolor=blue!60!black,
}

% Define colors for highlight boxes
\definecolor{findbg}{RGB}{235,245,255}
\definecolor{findbd}{RGB}{120,160,210}
\definecolor{warnbg}{RGB}{255,242,225}
\definecolor{warnbd}{RGB}{220,160,80}
\definecolor{tablehead}{RGB}{60,90,150}

\titleformat{\section}{\Large\bfseries\color{tablehead}}{\thesection}{1em}{}
\titleformat{\subsection}{\large\bfseries\color{tablehead!85}}{\thesubsection}{1em}{}

\newcommand{\codeword}[1]{\texttt{\small #1}}
\newcommand{\KEY}[1]{\textbf{\color{tablehead}#1}}

% Highlight boxes — simple colored vertical bar via \fcolorbox.
\newcommand{\finding}[2]{%
  \par\medskip
  \noindent\fcolorbox{findbd}{findbg}{%
    \begin{minipage}{\dimexpr\linewidth-2\fboxsep-2\fboxrule}
      \textbf{\color{findbd!40!black}#1.}\par\smallskip #2
    \end{minipage}}\par\medskip
}
\newcommand{\warningbox}[2]{%
  \par\medskip
  \noindent\fcolorbox{warnbd}{warnbg}{%
    \begin{minipage}{\dimexpr\linewidth-2\fboxsep-2\fboxrule}
      \textbf{\color{warnbd!40!black}#1.}\par\smallskip #2
    \end{minipage}}\par\medskip
}

\lstset{
  basicstyle=\ttfamily\footnotesize,
  breaklines=true,
  frame=single,
  framerule=0.3pt,
  rulecolor=\color{gray!50},
  backgroundcolor=\color{gray!5},
  xleftmargin=8pt,
  xrightmargin=8pt,
  aboveskip=8pt,
  belowskip=8pt,
}

\title{\Large\bfseries Cytokine Cascade Detection via Latent Space Geometry \\[0.3em]
\large Session Summary: Methods, Findings, and Open Challenges}
\author{Cytokine-MIL Project}
\date{19 May 2026}

\begin{document}

\maketitle

\begin{abstract}
\noindent
This report summarizes a working session on detecting cytokine signaling cascades
from a MIL-finetuned single-cell encoder. We document three complementary methods
investigated: (1) the PBS-RC \emph{directional bias} (\textbf{geo}) method for
candidate pair discovery, (2) the \textbf{cell-type ablation} method for
directionality determination --- found to be the strongest tool of the three ---
and (3) an alignment-based \textbf{inner-product / cosine} method for symmetric
pair detection. Diagnostic experiments on 8 random seeds showed that
\emph{pair detection itself is the current bottleneck}: of 11 pre-registered known
cascades, only 4 are recovered in any direction by the geo method at the 20\% top
threshold, and 0 in the forward direction at the 5\% threshold under per-source
top-3 filtering. The ablation method works reliably when handed a viable pair
list but cannot solve the upstream pair-selection problem. The session ends with
the implementation of an alignment-based detection variant designed to address
this bottleneck.
\end{abstract}

\section{Session Goal and Context}

The cytokine-MIL project trains an attention-based multiple-instance-learning
classifier on pseudo-tubes of PBMCs treated with one of 91 cytokines. The
hypothesis is that the trained model's latent geometry and dynamics encode
information about cytokine signaling \emph{cascades} --- cases where cytokine
$A$ induces a cellular state that subsequently produces cytokine $B$, with
some intermediate cell type $T$ acting as a \emph{relay}.

A cascade triple is therefore $(A, B, T)$: the source cytokine, the target
cytokine, and the relay cell type. The session focused on a two-stage cascade
detection pipeline already implemented (geo $\to$ ablation), diagnosed
multiple failure modes, applied iterative fixes, and finally designed an
alternative pair-detection method.

\paragraph{Validation ground truth.}
Eleven canonical cascade pairs (\codeword{KNOWN\_CASCADES} in
\codeword{scripts/analyze\_two\_stage\_results.py}) serve as the held-out
validation set:
\begin{quote}\small
IL-12 $\to$ IFN-$\gamma$, IL-1$\beta$ $\to$ IL-6, IL-2 $\to$ IL-15,
IL-33 $\to$ IL-13, IL-18 $\to$ IFN-$\gamma$, IL-21 $\to$ IL-10,
TNF-$\alpha$ $\to$ IL-6, IFN-$\alpha$1 $\to$ IFN-$\gamma$, IL-10 $\to$ IL-6,
IL-4 $\to$ IL-13, IL-27 $\to$ IFN-$\gamma$.
\end{quote}

\paragraph{Experiment grid.}
All experiments use 8 random seeds
(\codeword{42, 123, 7, 456, 789, 314, 271, 618}) with donors D1, D4--D12
for training and D2, D3 held out for validation. Each experiment trains
Stage 1 (encoder pre-training), Stage 2 (MIL with frozen encoder), and
Stage 3 (MIL with unfrozen encoder).

\section{Methods Investigated}

\subsection{Method 1: PBS-RC Directional Bias (``geo'')}

\paragraph{Goal.} Identify ordered candidate pairs $(A \to B)$ from the
latent space, with a candidate relay cell type $T$.

\paragraph{Definition.} For each tube of cytokine $C$, donor $d$, and
cell type $T$, the encoder produces a per-cell-type mean embedding
$\mu_{C,T,d} \in \mathbb{R}^{128}$. After subtracting the donor-pooled PBS
centroid for $T$ (PBS-RC), the directional bias of $A \to B$ via $T$
for donor $d$ is:
\[
b_{\mathrm{fwd}}^{(d)}(A \to B, T)
\;=\;
\bigl(\mu_{A,T,d} - \mu_A\bigr) \cdot \hat{u}_{A \to B},
\quad
\hat{u}_{A \to B} = \frac{\mu_B - \mu_A}{\|\mu_B - \mu_A\|}
\]
where $\mu_A$ is the pooled PBS-RC centroid of $A$ across cell types
(removing $A$'s generic effect, leaving only the $T$-specific cascade
component). The forward call uses a one-sided Wilcoxon signed-rank test
across donors with Bonferroni correction over cell types per ordered pair.

\paragraph{Encoder choice.} Initial diagnostic compared two encoders on
identical inputs:
\begin{itemize}
\item \textbf{Stage 1} encoder --- trained only on cell-type classification.
  Contains zero cytokine-specific information by construction.
\item \textbf{Stage 3} encoder --- MIL-finetuned. Cells are embedded in a
  cytokine-discriminative space.
\end{itemize}

\subsection{Method 2: Cell-Type Ablation (\textbf{the directional workhorse})}

\paragraph{Goal.} Given a candidate pair $(A, B)$, determine the direction
($A \to B$, $B \to A$, or shared) and identify the relay cell type $T^*$.

\paragraph{Definition.} For a tube treated with cytokine $A$ (target $B$),
compute the bag-level softmax probability $P(B \mid \mathrm{tube}_A)$ at
the full tube and at the tube with cell type $T$ ablated (all cells of
type $T$ removed). The relay score is:
\[
\mathrm{relay}(A \to B, T) \;=\; P(B \mid \mathrm{tube}_A^{\mathrm{full}})
\;-\; P(B \mid \mathrm{tube}_A^{\setminus T})
\]
A high positive value means cells of type $T$ in an $A$-treated tube are
\emph{necessary} for the model to predict $B$ --- direct functional
evidence of $A \to B$ via $T$.

The direction call compares max-relay across cell types in both directions
($A \to B$ vs $B \to A$). The pipeline does this for both forward and
reverse direction per (A, B) pair, so the ablation truly tests directionality
rather than assuming it from geo.

\finding{Why ablation is the strongest method we have}{
The ablation method has the property that \KEY{the model itself decides}
whether the cells of type $T$ in an $A$-treated tube carry the signal
for $B$. It is a direct functional test of the cascade hypothesis:
``does removing $T$ cells from $A$'s tube destroy the model's belief
in $B$?'' This sidesteps all assumptions about how the latent space
encodes cascades. As long as the model successfully learns to
discriminate cytokines, ablation gives a faithful direction signal.
The cost is computational: per pair $(A, B)$, each direction, each cell
type $T$, each donor, requires one forward pass --- $\sim$3.5 minutes
per ordered pair on the cluster (CPU, n=10 donors $\times$ n=18 cell
types $\times$ n=10 tubes per cytokine).
}

\subsection{Method 3: Alignment-Based Inner-Product / Cosine (new, designed this session)}

\paragraph{Goal.} A pair-detection method that does \emph{not} bake in a
directional assumption (symmetric scoring), addressing the pair-selection
bottleneck identified in this session.

\paragraph{Definition.} For each cytokine $C$, cell type $T$, donor $d$,
compute the PBS-RC centroid $v_{C,T,d} = \mu_{C,T,d} - \mu_{\mathrm{PBS},T}$.
For each unordered pair $\{A,B\}$, each $T$, each $d$:
\begin{align*}
\mathrm{cos}_d(A,B,T) &= \frac{\langle v_{A,T,d}, v_{B,T,d}\rangle}{\|v_{A,T,d}\|\cdot\|v_{B,T,d}\|}
\quad \text{(scale-free angular alignment)}\\
\mathrm{ip}_d(A,B,T)  &= \langle v_{A,T,d}, v_{B,T,d}\rangle
\quad \text{(magnitude-weighted)}
\end{align*}
Aggregate across donors by mean, then take max over $T$ for the
pair-level score. The relay $T^*(A,B)$ is the argmax cell type.

\paragraph{PCA-denoised variants.} The hypothesis is that of the 128
embedding dimensions, only a small fraction carry cascade-relevant
signal. To test this, we project the centroids to lower dimensions
(6 and 24) via global PCA on the stacked $\{v_{C,T,d}\}$ matrix, and
score in each subspace.

This produces six variants: \{cosine, inner product\} $\times$
\{6D, 24D, 128D\}.

\section{Diagnostic Results}

\subsection{Encoder choice: Stage 1 vs Stage 3 for geo detection}

The MIL-finetuned (Stage 3) encoder is substantially more informative than
the cell-type-only (Stage 1) encoder for cascade detection
(Table~\ref{tab:encoder-comparison}).

\begin{table}[h]
\centering
\caption{Stage 1 vs Stage 3 encoder geo recall on KNOWN\_CASCADES
(exp\_0\_seed42, all metrics on top-K\%). Stage 3 doubles
recall at every threshold under almost every metric.}
\label{tab:encoder-comparison}
\small
\begin{tabular}{lccccccc}
\toprule
\textbf{Metric} & \textbf{Stage} & \textbf{Recall@5\%} & & \textbf{Recall@10\%} & & \textbf{Recall@20\%} & \\
\midrule
max-W & S1 & 5/11 & & 5/11 & & 5/11 & \\
max-W & S3 & 4/11 & & 6/11 & & \textbf{8/11} & \\
\addlinespace
median-W & S1 & 0/11 & & 1/11 & & 2/11 & \\
median-W & S3 & 2/11 & & 4/11 & & \textbf{8/11} & \\
\addlinespace
mean-W & S1 & 1/11 & & 1/11 & & 4/11 & \\
mean-W & S3 & \textbf{4/11} & & \textbf{7/11} & & 7/11 & \\
\addlinespace
n\_ct\_at\_55 & S1 & 3/11 & & 5/11 & & 5/11 & \\
n\_ct\_at\_55 & S3 & 3/11 & & 6/11 & & 7/11 & \\
\bottomrule
\end{tabular}
\end{table}

\subsection{The W-saturation problem}

The original geo ranking used the \textbf{maximum} Wilcoxon W statistic
across cell types per pair. With $n=10$ donors, $W_{\max}=55$, and with
18 cell types, this maximum is reached for $\sim$14\% of pairs --- the
top of the ranking is saturated and provides no discrimination between
pairs.

\finding{Quantifying W-saturation}{
At Stage 3, of 8010 non-PBS ordered pairs:
\begin{itemize}
\item $W_{\max} = 55$ for \textbf{13.7\%} of pairs (1100 pairs all tied)
\item $W_{\max} \geq 50$ for 32.5\% of pairs
\item median-$W$ ranges 5--50, with much better spread (top 100 has 25 unique values)
\item mean-$W$ ranges 23--38, top 100 has 44 unique values
\end{itemize}
\textbf{Fix:} Switched ranking metric from max-$W$ to \textbf{mean-$W$
across cell types}. Mean-$W$ has the highest pair discrimination and
the best recall@5\% (4/11 known cascades).
}

\subsection{Per-source filtering: too aggressive at $K=3$}

To diversify the candidate set and ensure every cytokine gets
representation in the ablation budget, we tried \textbf{per-source top-K}
filtering: for each source cytokine $A$, keep its top-$K$ targets $B$ by
score. With $K=3$ and 90 cytokines, this yields $\sim$270 ordered pairs
per experiment (manageable for ablation).

\warningbox{Failure mode: known cascade pairs rank too low per source}{
For most known cascades, the cascade pair does not rank in the source
cytokine's top-3 (Table~\ref{tab:per-source-rank}). Stronger
``geometric outliers'' (e.g., IL-32-$\beta$ with its distinctive
B-cell-memory signature) dominate every source's top-3, even though
they are unlikely to be biological cascade partners. \KEY{0 of 11 known
cascades appear in the per-source top-3 forward direction list;
only 2 of 11 appear in any direction.}
}

\begin{table}[h]
\centering
\caption{Source cytokine's top-3 target list (exp\_0\_seed42, Stage 3,
mean-W ranking). For 9 of 11 known cascades, neither the source's
top-3 nor the target's top-3 includes the cascade partner.}
\label{tab:per-source-rank}
\footnotesize
\begin{tabular}{lll c}
\toprule
\textbf{Known} & \textbf{Source's top-3 targets} & \textbf{Target's top-3 sources} & \textbf{Found?} \\
\midrule
IL-12 $\to$ IFN-$\gamma$    & C5a, CT-1, TL1A             & IL-33, IL-8, IFN-$\alpha$1 & {\color{red}no} \\
IL-1$\beta$ $\to$ IL-6       & IFN-$\lambda$2, Leptin, FGF-$\beta$ & TNF-$\alpha$, IL-9, PRL    & {\color{red}no} \\
IL-2 $\to$ IL-15             & IFN-$\gamma$, IL-4, IFN-$\beta$ & IFN-$\beta$, IFN-$\gamma$, IFN-$\omega$ & {\color{red}no} \\
IL-33 $\to$ IL-13            & CD27L, IL-16, IFN-$\epsilon$ & IL-20, IL-27, IL-1$\alpha$ & {\color{red}no} \\
IL-18 $\to$ IFN-$\gamma$     & GITRL, CT-1, IFN-$\lambda$2 & IL-33, IL-8, IFN-$\alpha$1 & {\color{red}no} \\
IL-21 $\to$ IL-10            & IL-20, IL-22, ADSF          & TGF-$\beta$1, IFN-$\lambda$1, CD30L & {\color{red}no} \\
TNF-$\alpha$ $\to$ IL-6      & C3a, IL-31, IL-24           & \textbf{TNF-$\alpha$}, IL-9, PRL & {\color{green!50!black}reverse} \\
IFN-$\alpha$1 $\to$ IFN-$\gamma$ & M-CSF, RANKL, IL-17A    & IL-33, IL-8, \textbf{IFN-$\alpha$1} & {\color{green!50!black}reverse} \\
IL-10 $\to$ IL-6             & TGF-$\beta$1, IFN-$\lambda$1, CD30L & TNF-$\alpha$, IL-9, PRL & {\color{red}no} \\
IL-4 $\to$ IL-13             & IL-34, FLT3L, IFN-$\beta$   & IL-20, IL-27, IL-1$\alpha$ & {\color{red}no} \\
IL-27 $\to$ IFN-$\gamma$     & IL-3, FGF-$\beta$, IL-36Ra  & IL-33, IL-8, IFN-$\alpha$1 & {\color{red}no} \\
\bottomrule
\end{tabular}
\end{table}

\subsection{Ablation results on the per-source top-3 candidates}

Despite the geo selection missing most known cascades, the ablation step
on the chosen pairs identified several reproducible novel pairs
(Table~\ref{tab:ablation-strong-pairs}).

\begin{table}[h]
\centering
\caption{Pairs detected in $\geq$5 of 8 experiments via mean-W per-source
top-3 with ablation directionality. Numbers show experiment-level vote
counts ($A\to B$ : $B \to A$).}
\label{tab:ablation-strong-pairs}
\small
\begin{tabular}{lccl}
\toprule
\textbf{Pair} & \textbf{n\_exps} & \textbf{$A{\to}B$:$B{\to}A$} & \textbf{Biological interpretation} \\
\midrule
IL-9 $\to$ IL-10        & 6 & 5:1 & Treg / mast cell circuit \\
IL-23 $\to$ IFN-$\lambda$1 & 6 & 3:3 & Th17 $\to$ antiviral interferon \\
Noggin $\to$ IL-10      & 5 & 1:4 & BMP antagonist $\to$ anti-inflammatory \\
Decorin $\to$ TRAIL     & 5 & 3:2 & ECM proteoglycan $\to$ apoptosis \\
IL-15 $\to$ IL-1$\beta$ & 5 & 3:2 & NK activation $\to$ inflammasome \\
IL-36Ra $\to$ FasL      & 5 & 4:1 & IL-36 antagonist $\to$ apoptosis \\
IL-36Ra $\to$ IL-19     & 5 & 4:1 & IL-36 antagonist $\to$ IL-20 subfamily \\
IL-1$\alpha$ $\to$ CT-1 & 5 & 4:1 & Alarmin $\to$ cardiotrophin-1 \\
IL-1$\alpha$ $\to$ GM-CSF & 5 & 1:5 & Alarmin / myeloid expansion \\
GM-CSF $\to$ IFN-$\lambda$1 & 5 & 1:4 & Myeloid activation $\to$ IFN-$\lambda$ \\
GM-CSF $\to$ IL-27       & 5 & 2:3 & Myeloid activation $\to$ IL-27 \\
\bottomrule
\end{tabular}
\end{table}

\finding{The strongest novel discovery: IL-9 $\to$ IL-10}{
Detected in 6 of 8 experiments, with 5 of 6 informative votes calling
direction $A \to B$. IL-9 is a known driver of Treg expansion and IL-10
production --- this is a real immunological axis, validating that
ablation extracts genuine signal when handed reasonable candidates.
}

\section{Why Triplet Detection Is Hard: The Current Open Challenge}

The session's central diagnostic is that \emph{the cellular-relay
hypothesis is testable only after candidate triples $(A, B, T)$ are
proposed}. Three pair-selection strategies were tried:

\begin{table}[h]
\centering
\caption{Pair-selection strategy comparison.
Recall@5\% on KNOWN\_CASCADES; ``forward'' = $(A \to B)$ ordering
preserved; ``any'' = either direction matches.}
\label{tab:strategy-comparison}
\small
\begin{tabular}{lccc}
\toprule
\textbf{Strategy} & \textbf{Stage 3 max-W} & \textbf{Stage 3 mean-W} & \textbf{Stage 3 mean-W} \\
                  & \textbf{global top-5\%} & \textbf{global top-5\%} & \textbf{per-source top-3} \\
\midrule
Forward recall    & 4/11 & 4/11 & 0/11 \\
Any-direction recall & 4/11 & 4/11 & 2/11 \\
Pairs to ablate (per exp) & 400 & 400 & 270 \\
\bottomrule
\end{tabular}
\end{table}

\subsection{Why the geo method drops known cascades from its top}

The geo bias score measures \emph{how much} cells of type $T$ shift in the
direction of cytokine $B$ in an $A$-treated tube. Two known properties of
the dataset distort this score:
\begin{enumerate}[noitemsep]
\item A few cytokines (IL-32-$\beta$ above all) produce \emph{enormous}
  cell-type-specific shifts in cell-type embeddings. These dominate
  every source's top-K simply because they are geometric outliers.
\item Known cascade partners (e.g., IL-12 / IFN-$\gamma$, IL-1$\beta$ / IL-6)
  involve real but subtler shifts. Their pair score is real but
  middle-of-the-pack, so they get displaced by the outliers.
\end{enumerate}

\paragraph{Symmetry asymmetry.} Even when a known cascade has any
geometric signal, the geo direction is rarely the canonical one. For
instance, the dominant direction in the data is
``IFN-$\gamma$ $\to$ IL-12'' (because IFN-$\gamma$-treated cells shift
toward where IL-12 sits), not ``IL-12 $\to$ IFN-$\gamma$.'' This is a
property of \emph{latent space geometry}, not biology: the cytokine that
produces the bigger response also appears as the geometric ``source.''
Direction therefore should not be inferred from geo \emph{at all} ---
this is the ablation method's job (Section 2.2).

\subsection{Architecture of the proposed alignment method}

The inner-product / cosine method targets the pair-selection bottleneck
directly:
\begin{itemize}[noitemsep]
\item No directionality assumption (symmetric score).
\item Two metrics tested: cosine (angular) and raw inner product
  (magnitude-weighted).
\item Three dimensionalities tested: 128D (full), 24D (mild PCA), 6D
  (aggressive PCA --- testing the hypothesis that most dimensions are noise).
\item Output is fully compatible with the existing ablation pipeline:
  hand the top alignment pairs to ablation for directionality.
\end{itemize}

\warningbox{Submission status at session end}{
The alignment-based detection code is fully implemented, tested locally
(5/5 unit tests pass), committed to \codeword{main} (commit
\codeword{80f54bc}), and synced to the cluster. The SLURM submission
chain ran into intermittent OTP/auth failures during this session;
empirical results comparing the six variants against KNOWN\_CASCADES
are pending the next session's job execution.
}

\section{Implementation Summary}

\subsection{Codebase changes from this session}

\begin{table}[h]
\centering
\caption{Files added or modified.}
\label{tab:files}
\footnotesize
\begin{tabular}{p{6.5cm}p{8cm}}
\toprule
\textbf{File} & \textbf{Purpose} \\
\midrule
\codeword{cytokine\_mil/analysis/encoder\_cache.py} \textit{(new)}
& Shared helpers: \codeword{build\_cache}, \codeword{compute\_pbs\_centroids},
  \codeword{\_embed\_tube}, \codeword{\_load\_label\_encoder}. Extracted from
  \codeword{run\_geo\_extract.py} for reuse. \\
\codeword{cytokine\_mil/analysis/pair\_alignment.py} \textit{(new)}
& Core alignment computation: per-(C,T,d) PBS-RC centroids, PCA projections,
  pair scoring under cosine and inner-product metrics, top-pairs formatting. \\
\codeword{scripts/run\_geo\_extract.py} \textit{(refactored)}
& Now imports cache helpers from \codeword{encoder\_cache.py}. \\
\codeword{scripts/run\_inner\_product\_pairs.py} \textit{(new)}
& Per-experiment alignment runner. Six variant outputs per experiment. \\
\codeword{scripts/analyze\_inner\_product\_results.py} \textit{(new)}
& Cross-experiment recall@K analysis, dim comparison plots. \\
\codeword{scripts/filter\_top\_pairs.py} \textit{(modified)}
& Switched ranking from max-W to mean-W across cell types. \\
\codeword{scripts/analyze\_two\_stage\_results.py} \textit{(modified)}
& Threshold sweep now uses mean-W (with fallbacks). \\
\codeword{slurm/inner\_product/01\_run.slurm} \textit{(new)}
& Array job 0--7, 32 GB, 4 CPUs, 1 h. \\
\codeword{slurm/inner\_product/02\_analyze.slurm} \textit{(new)}
& Cross-experiment analysis, 8 GB, 30 min. \\
\codeword{slurm/inner\_product/meta\_submit.slurm} \textit{(new)}
& Chains 01 $\to$ 02 via \codeword{sbatch --dependency=afterok}. \\
\codeword{tests/test\_pair\_alignment.py} \textit{(new)}
& 5 unit tests: shape, symmetry, identity-cosine, PCA dims, schema. \\
\bottomrule
\end{tabular}
\end{table}

\subsection{Reusable utilities}

\begin{table}[h]
\centering
\caption{Key utilities, with function signatures, that future work
should reuse rather than reimplement.}
\label{tab:utilities}
\footnotesize
\begin{tabular}{p{5cm}p{10cm}}
\toprule
\textbf{Location} & \textbf{Signature} \\
\midrule
\codeword{encoder\_cache.py} &
\codeword{compute\_pbs\_centroids(encoder, manifest, gene\_names, device)}
$\to$ \codeword{Dict[cell\_type, ndarray]} \\
\codeword{encoder\_cache.py} &
\codeword{build\_cache(encoder, manifest, gene\_names, label\_encoder, device)}
$\to$ \codeword{list[dict]} (each tube as \codeword{\{H, label, cell\_types, donor\}}) \\
\codeword{pair\_alignment.py} &
\codeword{compute\_per\_atd\_centroids(cache, label\_encoder, pbs\_ct\_means, train\_donors)}
$\to$ \codeword{Dict[(C,T,d), ndarray]} \\
\codeword{pair\_alignment.py} &
\codeword{fit\_pca\_projections(centroids, n\_components\_list)}
$\to$ \codeword{Dict[dim, Dict[(C,T,d), ndarray]]} \\
\codeword{pair\_alignment.py} &
\codeword{compute\_pair\_scores(centroids, metrics)}
$\to$ \codeword{(pair\_scores, relay\_T, full\_table)} \\
\codeword{latent\_geometry.py} &
\codeword{compute\_directional\_bias\_per\_donor(cache, le, pbs\_ct, train\_donors, direction\_mode)} \\
\codeword{latent\_geometry.py} &
\codeword{test\_directional\_significance(bias\_result, le, alpha)} \\
\codeword{analyze\_cell\_type\_ablation.py} &
Ablation entry point; consumes \codeword{top\_pairs.json} schema $\{A, B, relay\_cell\_type, \ldots\}$. \\
\bottomrule
\end{tabular}
\end{table}

\section{Conclusions and Forward Plan}

\subsection{What works}

\begin{enumerate}[label=\textbf{C\arabic*.},leftmargin=*]
\item \KEY{Cell-type ablation is the most reliable component of the
pipeline.} Given a viable pair, ablation tests directionality directly
via the model's own predictions. It is computationally expensive but
biologically faithful. \textbf{Future work should treat ablation as the
``ground truth'' direction call} and focus on improving upstream pair
selection.

\item The Stage-3 MIL-finetuned encoder doubles known-cascade recall
over the Stage-1 encoder. Future work should use \codeword{model\_stage3.pt}.

\item Mean-W across cell types is a far better pair ranking metric
than max-W (which saturates) or median-W (which has poor pair-level
spread).

\item The pipeline is now properly modularized: \codeword{encoder\_cache.py}
holds the shared cache logic, so new pair-detection methods only need
to implement scoring on the centroid dictionary.
\end{enumerate}

\subsection{What does not work yet}

\begin{enumerate}[label=\textbf{X\arabic*.},leftmargin=*]
\item \KEY{Pair selection from the geo bias score misses known cascades.}
Per-source top-3 yields 0/11 forward recall; global top-5\% yields 4/11
in any direction. The latent geometry contains too many strong outlier
cytokines (notably IL-32-$\beta$) that dominate the rankings.

\item \KEY{The cellular-relay triplet $(A, B, T)$ is the right
biological unit, but our current scores are pair-level and only weakly
$T$-conditional.} The relay cell type is recovered as an argmax, with
no calibration on its plausibility.

\item Direction information from geo is unreliable. Asymmetry between
$(A \to B)$ and $(B \to A)$ geo scores reflects which cytokine has the
bigger response, not the direction of the cascade.
\end{enumerate}

\subsection{Specific next actions for the future session}

\begin{enumerate}[label=\textbf{N\arabic*.},leftmargin=*,itemsep=2pt]
\item Run the alignment-based detection pipeline submitted at session
end (\codeword{slurm/inner\_product/meta\_submit.slurm}). Compare the
six (metric, dim) variants against KNOWN\_CASCADES via
\codeword{analyze\_inner\_product\_results.py} $\to$ \codeword{recall\_curves.png}.
\item If a variant achieves recall@5\% $\geq 4/11$, hand its
\codeword{top\_pairs\_*.json} to the existing ablation pipeline
(\codeword{scripts/analyze\_cell\_type\_ablation.py}) for direction
calling.
\item Diagnose IL-32-$\beta$ specifically: investigate whether it is
a real biological anchor or a training-set artifact (donor-specific?
batch effect? unusual transcriptional signature in the encoder).
\item Consider running ablation on a deliberately constructed pair set
containing the 11 known cascades, to establish a clean ``ablation
recovers direction'' benchmark independent of pair selection.
\end{enumerate}

\section*{Appendix: Key Quantitative Snapshots}

\begin{table}[h]
\centering
\caption{Quantitative summary table for cross-method comparison
(Stage 3 encoder, exp\_0\_seed42).}
\label{tab:method-summary}
\small
\begin{tabular}{lcccc}
\toprule
\textbf{Method} & \textbf{Recall@5\%} & \textbf{Score range} &
\textbf{Discrimination} & \textbf{Direction call} \\
\midrule
geo max-W                & 4/11 (any dir) & saturates at 55 & poor    & implicit, unreliable \\
geo median-W             & 2/11           & 5--50          & moderate & ditto \\
geo mean-W               & 4/11           & 23--38         & good    & ditto \\
per-source top-3 (mean-W) & 0/11 (fwd)   & ditto          & --       & dominated by outliers \\
ablation only             & --            & --             & --       & \KEY{direct, reliable} \\
inner-product (designed)  & TBD           & TBD           & TBD     & \KEY{none (symmetric)} \\
\bottomrule
\end{tabular}
\end{table}

\begin{table}[h]
\centering
\caption{The eight experiments and their seeds.}
\label{tab:experiments}
\small
\begin{tabular}{cc}
\toprule
\textbf{Exp index} & \textbf{Seed} \\
\midrule
0 & 42 \\
1 & 123 \\
2 & 7 \\
3 & 456 \\
4 & 789 \\
5 & 314 \\
6 & 271 \\
7 & 618 \\
\bottomrule
\end{tabular}
\end{table}

\paragraph{Training donors.} D1, D4, D5, D6, D7, D8, D9, D10, D11, D12 (10 donors).
\paragraph{Validation donors (held out).} D2, D3 (2 donors).

\end{document}
